\section{Discussion}
\label{sec:discussion}

\subsection{Why does \COT help here when ``curse of CoT'' says it should not?}

The mechanism is task shape, not task novelty. \addrtask is not rule
induction. The grid is given verbatim and the question asks for a
verifiable property of that grid; the model does not have to infer a
transformation. For \qpoint, \qlookup, and \qneighbor the model can
index into the prompt directly, and \Direct prompting succeeds at
$\ge 95\%$ accuracy. For \qcount and \qrowmax the model has to scan,
accumulate, and aggregate---exactly the operations that
\citet{sprague2025tocot} identify as the regime where CoT actually buys
accuracy. The fact that \Direct \qcount errors are almost all
off-by-one is the textbook failure mode that an externalised reasoning
trace was designed to fix.

This refines, rather than contradicts, the curse-of-CoT story
\citep{zheng2025curse}. The \miniarc control reproduces
$\Direct{>}\COT$ in the rule-induction regime where the curse applies;
the \addrtask result shows the ordering inverts as soon as the query
becomes a scan over the grid. The implication for ICL benchmarking is
that pooling rule-induction and addressing queries into a single
``pattern-induction'' bucket conflates two very different effects.

\subsection{Repetition does not bridge the aggregation gap}

\Tabref{tab:qtype} reveals that even at $K{=}8$ same-grid demos,
\gpt \Direct only reaches 0.44 on \qcount and \haiku \Direct reaches
0.875. The demonstration pool deliberately includes other \qcount
examples about the same grid with correct answers, but the model
still cannot generalise the counting procedure to a new token within
the same grid. Repetition lifts \emph{base accuracy} on the random-access
queries, but it does not install a counting procedure---an observation
consistent with the surface self-reinforcement view of
\citet{yan2024repetitions} but at odds with the strong reading of
\citet{liu2025sparks}, which would predict that demonstrations bind a
symbolic understanding of the grid that transfers across queries.

\subsection{Practical takeaway}

For workflows that ground LLM decisions in small text grids (UI dumps,
board games, ARC-style puzzles, tabular layouts), enabling \COT is a
strictly cheaper reliability lever than padding the prompt with
identical-grid demonstrations: it closes the aggregation gap that
demonstrations cannot, and the token-cost penalty is dominated by the
accuracy gain at any realistic price ratio (\figref{fig:cost}).

\subsection{Limitations}

\para{Two models only.}
Both \gpt and \haiku are relatively small, ``fast'' production models.
Larger reasoning models (o-series, Sonnet 4.5 with extended thinking,
Gemini 2.5 Pro) plausibly close more of the gap even without
explicit \COT; the addressing gap might shrink as base capacity grows.

\para{Single grid size.}
We evaluate at $7{\times}7$ only. \citet{bai2025matrix} report a
42--84\% accuracy drop as grids grow; our results may be optimistic for
larger grids, particularly for \qcount and \qrowmax where the scan
length grows quadratically.

\para{Token alphabet is colour-coded.}
The alphabet $\{R, B, G, Y\}$ carries colour semantics in pre-training.
Truly arbitrary symbols (non-ASCII glyphs, random Unicode) may degrade
either condition.

\para{Strict exact-match scoring.}
We parse and score only the model's \emph{final answer}, even when \COT
intermediate reasoning matches the gold. A well-reasoned response that
fumbles the final-answer line counts as wrong; this likely slightly
under-states \COT accuracy.

\para{Power for fine-grained contrasts.}
With $n{=}80$ paired observations per cell, contrasts at the ceiling
(\eg \COT $K{=}1$ vs.\ $K{=}2$) are noisy: a single flipped item can
move the estimate by 1.25 percentage points.

\para{\miniarc compliance.}
\haiku ignored the ``\Direct only, no explanation'' instruction on
\miniarc and produced reasoning prose our grid parser could not
extract. The \miniarc numbers should therefore be read as
direction-of-effect, not as a benchmark. The \addrtask format had no
such compliance issue.

\para{No vision modality.}
All grids are text. The perceptual encoding story behind the
forgotten-polygons line~\citep{rudman2025polygons} does not apply
directly to our setting.
